Using the notion of weak derivatives, we can weaken the requirements for solutions to a PDE. Consider, for instance, the heat equation in $(0, T) \times \Omega$, where $\Omega \subseteq \R^m$ is open and bounded with a smooth boundary, with homogeneous Neumann boundary conditions,

\begin{equation}
\label{eq:strong_heat}
\begin{aligned}
    u_t &= \Delta u         & \text{in } &(0, T) \times \Omega,\\
    \pderiv[u]{\nu} &= 0    & \text{on } &[0, T) \times \partial \Omega.
\end{aligned}
\end{equation}

\noindent
For a function $u$ to be a classical solution, we require $u$ to be of class $C^{1,2}((0, T) \times \Omega) \cap C^{0}([0, T] \times \overline{\Omega})$, i.e. continuously differentiable with respect to $t$, twice continuously differentiable with respect to $x$ and continuously extendable to $[0, T] \times \overline{\Omega}$. We can multiply the heat equation by a smooth \todo{smooth, not test function right?}function $\phi \in C^{\infty}(\Omega)$, integrate over $\Omega$ and integrate by parts to derive

\begin{equation*}
    \int_{\Omega} u_{t} \phi = \int_{\Omega} \Delta u \phi = \int_{\partial \Omega} \frac{\partial u}{\partial \nu} \phi \d S - \int_{\Omega} \nabla u \cdot \nabla \phi = - \int_{\Omega} \nabla u \cdot \nabla \phi .
\end{equation*}

\noindent
We can also write this as 

\begin{equation}\label{eq:weak_heat}
    \inner{u_t, \phi}_{\Ltwo{\Omega}} = -\inner{\nabla u, \nabla \phi}_{\Ltwo{\Omega; \R^d}} ,
\end{equation}

\noindent
and by \cref{prop:approx_by_smooth}, we can show that this should also hold for every $\phi \in H^1(\Omega)$. Now note that for \cref{eq:weak_heat} to be well-defined, we don't need $u$ to be twice continuously differentiable anymore. In fact, we don't even need $\nabla u$ to be continuous. Instead, square-integrability of both $\nabla u$ and $u_t$ is sufficient. This means we can substantially weaken the requirements for $u$, and instead ask if there exists a function $u \in H^1(0, T; H^1(\Omega))$.